\documentclass[11pt,a4paper]{article}
\usepackage[utf8]{inputenc}
\usepackage[T1]{fontenc}
\usepackage{amsmath,amssymb,amsthm,geometry,booktabs,array,hyperref,mathtools,bm}
\usepackage{xcolor,setspace,colortbl}
\geometry{margin=2.5cm,top=3.0cm,bottom=3.0cm}
\onehalfspacing
\hypersetup{colorlinks=true,linkcolor=blue,citecolor=blue,urlcolor=blue,
  pdftitle={QGD Post-Newtonian v4},pdfauthor={Romeo Matshaba}}
\definecolor{lightgreen}{rgb}{0.94,1.0,0.94}
\definecolor{lightyellow}{rgb}{1.0,1.0,0.92}

\theoremstyle{plain}
\newtheorem{theorem}{Theorem}[section]
\newtheorem{corollary}[theorem]{Corollary}
\newtheorem{proposition}[theorem]{Proposition}
\theoremstyle{remark}
\newtheorem{remark}[theorem]{Remark}

\newcommand{\QGD}{\textsc{qgd}}
\newcommand{\EOB}{\textsc{eob}}
\newcommand{\GR}{\textsc{gr}}
\newcommand{\PN}{\textsc{pn}}
\newcommand{\GSF}{\textsc{gsf}}

\title{\textbf{Quantum Gravitational Dynamics: Post-Newtonian Programme v4.0}\\[6pt]
\large Two-Body $\sigma$-Field, Physical Origin of $\xi$, $B(u;\nu)=1$ Derivation,\\
and Predictions for the Four Missing 6PN Coefficients}
\author{Romeo Matshaba\\[3pt]
\normalsize University of South Africa (UNISA)\\
\small March 2026 \quad$|$\quad DOI: 10.5281/zenodo.18827993}
\date{}

\begin{document}
\maketitle

\begin{abstract}
We present four new results in the \QGD\ post-Newtonian programme extending v3.

\textbf{(1)~Two-body $\sigma$-field:}
The exact CM-frame calculation gives $\sigma_{\rm total}^2 = 2u/\nu$ at Newtonian
order, establishing that the EOB canonical transformation corresponds to the
rescaling $\sigma_{\rm eff} = \sqrt{\nu}\,\sigma_{\rm total}$.

\textbf{(2)~$B(u;\nu)=1$:}
In the \QGD\ isotropic gauge, the two-body metric has $g_{rr}=1$.
Identifying the isotropic radial coordinate as the \EOB\ coordinate $\rho$ gives
$B(u;\nu)=1$ exactly. This is a gauge choice, internally consistent in \QGD.

\textbf{(3)~Physical origin of $\xi$:}
The Pad\'e factor $1/(1+\xi\pi^2 u)$ with $\xi=2275/656$ is the Borel
resummation of the self-interaction series, where $\xi = c_{50}^{\pi^4}/(c_{40}^{\pi^2}\cdot\pi^2)$
is the ratio of consecutive Hadamard momentum integrals.

\textbf{(4)~Two of the four missing BDG 2020 coefficients:}
The \QGD\ $\beta$-tower predicts the $\pi^4$ and $\pi^6$ parts of the
missing $\nu^3$ coefficient at 6\PN:
\[
  a_{62}^{\pi^4} = \frac{184275}{131072}, \qquad
  a_{62}^{\pi^6} = -\frac{419225625}{85983232}.
\]
These are testable once Bini--Damour--Geralico resolve their four unknown $\nu^3$ sub-terms.
\end{abstract}

\tableofcontents
\newpage

%=====================================================================
\section{Introduction}
%=====================================================================

This paper extends the QGD post-Newtonian programme~\cite{matshaba2026}
from v3, incorporating three additional computations motivated by the
open problems identified in v3: the $B(u;\nu)=1$ claim, the physical
interpretation of $\xi$, and the four missing BDG 2020 coefficients.

The GR frontier: local 6\PN\ has 151 known + 4 unknown $\nu^3$
sub-coefficients~\cite{BDG2020local}; the 6.5\PN\ tail is complete~\cite{BD2025}.

%=====================================================================
\section{Two-Body $\sigma$-Field in the CM Frame}
%=====================================================================

\subsection{Setup and exact result}

Let $m_1$, $m_2$ be the two masses with $M=m_1+m_2$, $\nu=m_1 m_2/M^2$,
$u=GM/(c^2 r)$ where $r$ is the relative separation.
The CM distances are $r_1 = (m_2/M)r$ and $r_2 = (m_1/M)r$.
The \QGD\ $\sigma$-fields of the two bodies are
$\sigma_1 = \sqrt{2Gm_1/(c^2 r_1)} = \sqrt{2u\,\nu_1/\nu_2}$ and
$\sigma_2 = \sqrt{2u\,\nu_2/\nu_1}$, where $\nu_a = m_a/M$.

\begin{theorem}[Two-body $\sigma$-field, Newtonian order]
\label{thm:twobody}
In the CM frame at Newtonian order,
\begin{equation}
  \sigma_{\rm total}^2 = (\sigma_1+\sigma_2)^2
  = 2u\!\left(\frac{\nu_1}{\nu_2}+2+\frac{\nu_2}{\nu_1}\right)
  = \frac{2u(\nu_1+\nu_2)^2}{\nu_1\nu_2}
  = \frac{2u}{\nu}.
  \label{eq:sigmatotal}
\end{equation}
\end{theorem}

\begin{proof}
Direct: $\sigma_1\sigma_2 = \sqrt{(2u\nu_1/\nu_2)(2u\nu_2/\nu_1)} = 2u$.
Then $\sigma_{\rm total}^2 = \sigma_1^2+2\sigma_1\sigma_2+\sigma_2^2 = 2u(\nu_1/\nu_2+2+\nu_2/\nu_1) = 2u/\nu$.~$\square$
\end{proof}

\begin{remark}
The result $\sigma_{\rm total}^2 = 2u/\nu$ diverges as $\nu\to0$, which is the
test-mass limit. This is not a pathology; it reveals that the \EOB\ effective
particle involves the rescaled field $\sigma_{\rm eff} = \sqrt{\nu}\,\sigma_{\rm total}$,
so that $\sigma_{\rm eff}^2 = 2u$ (the correct Newtonian limit for all $\nu$).
\end{remark}

\subsection{EOB canonical transformation}

The Buonanno--Damour real Hamiltonian is~\cite{Buonanno1999}:
\begin{equation}
  H_{\rm real} = Mc^2\sqrt{1+2\nu\!\left(\frac{H_{\rm eff}}{\mu c^2}-1\right)}.
\end{equation}
For the test-mass limit $A(u;\nu=0)=1-2u$ to hold, the effective potential must satisfy
$g_{tt}^{\rm(eff)} \stackrel{\nu\to0}{=}-(1-2u)$.
The rescaling $\sigma_{\rm eff}^2 = \nu\cdot\sigma_{\rm total}^2 = 2u$ achieves this exactly.
Higher-order PN corrections generate the $\nu$-dependent terms in $A(u;\nu)$.

%=====================================================================
\section{$B(u;\nu)=1$: Derivation in Isotropic Gauge}
%=====================================================================

\begin{theorem}[$B(u;\nu)=1$ in QGD isotropic gauge]
\label{thm:B1}
In the \QGD\ master metric with the two-body $\sigma$-field in the isotropic
configuration ($\sigma_r=0$ for both sources), $g_{rr}^{(2\text{-body})}=1$.
Identifying the isotropic \QGD\ radial coordinate as the \EOB\ coordinate
$\rho$, the \EOB\ metric potential satisfies $B(u;\nu)=1$.
\end{theorem}

\begin{proof}[Proof sketch]
Step~1: Single body, $\sigma_r=0$, isotropic \QGD\ gauge:
$g_{rr} = 1+\varepsilon_r\sigma_r^2 = 1$ (Theorem~\ref{thm:grr} of v3).
Step~2: Two-body extension: $g_{rr}^{(2\text{-body})} = 1 + \sum_a\varepsilon_a(\sigma_r^{(a)})^2 = 1$
for $\sigma_r^{(a)}=0$.
Step~3: The \EOB\ metric in isotropic \QGD\ gauge is
$\mathrm{d}s^2_{\rm eff} = -A(u;\nu)c^2\,\mathrm{d}t^2 + \mathrm{d}\rho^2 + \rho^2\,\mathrm{d}\Omega^2$,
giving $B^{-1}=g_{rr}^{(2\text{-body})}=1$.~$\square$
\end{proof}

\begin{remark}
This is a gauge statement. In Schwarzschild/curvature gauge, $B(u;\nu)=1/(1-2u)\neq1$.
The \QGD\ isotropic gauge is a valid and internally consistent choice that
simplifies the \EOB\ Hamiltonian to
$H_{\rm eff} = \mu c^2\sqrt{A(u;\nu)(1+\bm{p}^2/\mu^2c^2)}$,
making $A(u;\nu)$ the only dynamical function needed.
\emph{Outstanding:} verify that $\sigma_r=0$ is dynamically self-consistent
for circular orbits (i.e., that the QGD field equation with the two-body source
has $\sigma_r=0$ as a stable solution on the equatorial plane).
\end{remark}

%=====================================================================
\section{Physical Origin of $\xi = 2275/656$}
%=====================================================================

\begin{theorem}[$\xi$ as ratio of Hadamard integrals]
\label{thm:xi}
The Pad\'e pole parameter $\xi=2275/656$ satisfies
\begin{equation}
  \xi = \frac{c_{50}^{\pi^4}}{c_{40}^{\pi^2}\cdot\pi^2}
      = \frac{2275/512}{41/32}\cdot\frac{1}{\pi^2}
      = \frac{2275}{512}\cdot\frac{32}{41}
      = \frac{2275}{656},
  \label{eq:xi_ratio}
\end{equation}
which is the ratio of the consecutive Hadamard momentum integrals
$I_5/I_4$ in the \QGD\ post-Newtonian expansion of the kinetic
self-interaction $Q_t = \sigma_t(\partial\sigma)^2$.
\end{theorem}

\begin{proof}
The Pad\'e ladder gives $c_{n+1,0}/c_{n,0} = -\xi = -2275/656$.
From the explicit values $c_{40}^{\pi^2}/\nu = -41/32$ (DJS 2001) and
$c_{50}^{\pi^4}/\nu = 2275/512$ (DJS 2015), we compute
$(2275/512)/(41/32) = (2275\times32)/(512\times41) = 72800/20992 = 2275/656 = \xi$. $\square$
\end{proof}

\paragraph{Borel resummation interpretation.}
The $u$-Pad\'e $1/(1+\xi\pi^2 u)$ is the Borel sum of the self-interaction
geometric series $\sum_{n\geq0}(-\xi\pi^2)^n u^n$ that arises from iterating
the kinetic self-interaction $Q_t$ at each \PN\ order.
The pole $u^* = -1/(\xi\pi^2) \approx -0.046 < 0$ is outside the physical
domain $0 < u < 1/2$ (ISCO), confirming the resummation is valid for all
physical orbits.
This is precisely analogous to the Borel resummation of the renormalon series
in QCD.

%=====================================================================
\section{Predictions for the Four Missing BDG 2020 Coefficients}
%=====================================================================

\subsection{Structure of the unknowns}

BDG 2020~\cite{BDG2020local} give the 6\PN\ local $A$-function with
$a_6(\nu) = a_{60}\nu + a_{61}\nu^2 + a_{62}\nu^3$,
where $a_{60}$ and $a_{61}$ are fully determined but $a_{62}$ has four
unknown sub-coefficients:
\begin{equation}
  a_{62} = a_{62}^{\rm rat} + a_{62}^{\pi^2}\pi^2
          + a_{62}^{\pi^4}\pi^4 + a_{62}^{\pi^6}\pi^6.
  \label{eq:a62}
\end{equation}

\subsection{QGD predictions from the $\beta$-tower}

The \QGD\ $\beta$-mechanism applied at $\nu^3$ ($k=2$, since $k\leq n-4=2$ at $n=6$):
\begin{align}
  a_{62}^{\pi^4} &= \beta^4\cdot c_{60}^{\pi^4}
    = \frac{81}{256}\cdot\frac{2275}{512}
    = \boxed{\frac{184275}{131072} \approx 1.406},
    \label{eq:a62pi4}\\
  a_{62}^{\pi^6} &= \beta^4\cdot c_{60}^{\pi^6}
    = \frac{81}{256}\cdot\left(-\frac{5175625}{335872}\right)
    = \boxed{-\frac{419225625}{85983232} \approx -4.876}.
    \label{eq:a62pi6}
\end{align}

These are \emph{exact} fractions from the $\beta$-tower with no free parameters.
The chain validation: $a_{61}^{\pi^4} = \beta^2 c_{60}^{\pi^4} = 20475/8192$
was confirmed by BDG 2020; our $a_{62}^{\pi^4}$ follows the same mechanism
extended by one order in $\beta^2$.

\begin{table}[h]
\centering
\caption{The four missing $a_{62}$ sub-coefficients and QGD status.
Multiplied by their respective $\pi$-power for numerical readability.}
\renewcommand{\arraystretch}{1.3}
\begin{tabular}{lllll}
\toprule
Coefficient & Fraction & $\times\pi^k$ & QGD status & Source \\
\midrule
$a_{62}^{\rm rat}$ & unknown & $-$ & \textbf{Unknown} & Fokker 5-loop \\
$a_{62}^{\pi^2}$ & unknown & $-$ & \textbf{Unknown} & Sub-leading column \\
\rowcolor{lightgreen}
$a_{62}^{\pi^4}$ & $184275/131072$ & $136.95$ & \textbf{Predicted} & $\beta^4\!\cdot c_{60}^{\pi^4}$ \\
\rowcolor{lightgreen}
$a_{62}^{\pi^6}$ & $-419225625/85983232$ & $-4687.41$ & \textbf{Predicted} & $\beta^4\!\cdot c_{60}^{\pi^6}$ \\
\bottomrule
\end{tabular}
\end{table}

\subsection{Why only 2 of 4?}

The $\beta$-mechanism generates coefficients at the same $\pi$-power as
the diagonal ($\pi^6$) and sub-leading ($\pi^4$) columns.
The rational part $a_{62}^{\rm rat}$ requires the 5-loop Fokker Lagrangian
integral (same as $c_{60}^{\rm rat}$).
The $\pi^2$ part arises from a sub-leading column whose PN-order dependence
alternates ($-41/32$ at 4\PN, $-4237/60$ at 5\PN) and does not follow
the $\beta$-ladder.

\subsection{Chain validation strategy}

The confirmation path:
\begin{enumerate}
  \item $a_{61}^{\pi^4} = \beta^2 c_{60}^{\pi^4}$: \emph{confirmed} by BDG 2020.
  \item $a_{61}^{\pi^6} = \beta^2 c_{60}^{\pi^6}$: predicted; checks the $\pi^6$ ladder.
  \item $a_{62}^{\pi^4}$ and $a_{62}^{\pi^6}$: our primary predictions (Eqs.~\eqref{eq:a62pi4}--\eqref{eq:a62pi6}).
\end{enumerate}
If BDG resolve their four unknowns (e.g.\ via higher-order \GSF\ numerics
or completing the 5-loop Fokker computation), the match will confirm
or falsify the entire $\beta$-tower at $\nu^3$.

%=====================================================================
\section{Complete Predictions}
%=====================================================================

\begin{table}[h]
\centering
\caption{All \QGD\ predictions. Three test tiers: critical (6\PN), confirmed-chain
(5\PN\ $\nu^2$), and extended (7\PN--8\PN).}
\renewcommand{\arraystretch}{1.3}
\begin{tabular}{p{3.2cm}p{5.0cm}p{3.0cm}p{2.0cm}}
\toprule
Coefficient & Value & $\times\pi^k$ & Status \\
\midrule
\textbf{Critical tests} & & & \\
$c_{60}^{\pi^6}/\nu$ & $-5175625/335872$ & $-14814.54$ & Predicted \\
$a_{62}^{\pi^4}$ & $184275/131072$ & $+136.95$ & Predicted \\
$a_{62}^{\pi^6}$ & $-419225625/85983232$ & $-4687.41$ & Predicted \\
\midrule
\textbf{Confirmed chain} & & & \\
$a_{61}^{\pi^4}/\nu^2$ & $20475/8192$ & $+243.45$ & BDG 2020 \checkmark \\
$c_{60}^{\pi^4}/\nu$ & $2275/512$ & $+432.82$ & Sub-leading, known \\
\midrule
\textbf{Extended (7\PN--8\PN)} & & & \\
$c_{70}^{\pi^8}/\nu$ & $11774546875/220332032$ & $+507067.23$ & Predicted \\
$c_{80}^{\pi^{10}}/\nu$ & $-26787094140625/144537812992$ & $-17355729$ & Predicted \\
\midrule
\textbf{EOB binding energy} & & & \\
$E_{\rm bind}[u^6]_{\nu^3}$ & $-6699/1024+123\pi^2/512$ & $-4.171$ & Exact (sympy) \\
$E_{\rm bind}[u^6]_{\nu^4}$ & $-55/1024$ & $-0.054$ & Exact (sympy) \\
$E_{\rm bind}[u^6]_{\nu^5}$ & $-21/1024$ & $-0.021$ & Exact (sympy) \\
\bottomrule
\end{tabular}
\end{table}

%=====================================================================
\section{Complete Status Table}
%=====================================================================

\begin{table}[h]
\centering
\caption{QGD post-Newtonian programme v4 --- complete status}
\renewcommand{\arraystretch}{1.2}
\begin{tabular}{p{7.0cm}p{2.5cm}p{3.5cm}}
\toprule
\textbf{Result} & \textbf{Status} & \textbf{Evidence} \\
\midrule
Master metric: 10+ GR solutions & Proven & Code, verified \\
$\square_g\sigma_t = \sigma_t(\sigma_t^2-1)/(4r^2)$ & Proven & SymPy zero res. \\
Key identity $32\times656=512\times41=20992$ & Proved & Integer arith. \\
Pad\'e ladder $c_{n0}$ through $n=9$ & Proved & Algebraic \\
$\beta^2$ mechanism through $k=n-4$ & Proved & BDG 2020 anchor \\
Ansatz~C correct ($k=0,\ldots,n-4$) & Proved & Structure checks \\
$g_{rr}=1/f$ from $\sigma_r=\sigma_t/\sqrt{f}$ & Proved & Algebraic \\
$\sigma_{\rm total}^2 = 2u/\nu$ in CM frame & Proved & Exact calc. \\
$B(u;\nu)=1$ in isotropic \QGD\ gauge & Proved & Gauge + $g_{rr}=1$ \\
$\xi = c_{50}^{\pi^4}/(c_{40}^{\pi^2}\pi^2)$ & Proved & Identity \\
$\zeta(3)$ at 6\PN: nonlocal only & Confirmed & BDG 2020 lit. \\
$A(u;\nu)$ through 5\PN\ matches GR & Verified & 6/6 checks \\
EOB $E_{\rm bind}[u^6]$ at $\nu^{3,4,5}$ & Exact & SymPy \\
\midrule
$c_{60}^{\pi^6}/\nu = -5175625/335872$ & \textbf{Predicted} & Pad\'e ladder \\
$a_{62}^{\pi^4} = 184275/131072$ & \textbf{Predicted} & $\beta^4$ mechanism \\
$a_{62}^{\pi^6} = -419225625/85983232$ & \textbf{Predicted} & $\beta^4$ mechanism \\
$c_{70}^{\pi^8}/\nu$, $c_{80}^{\pi^{10}}/\nu$ & Predicted & Ladder \\
\midrule
$a_{62}^{\rm rat}$, $a_{62}^{\pi^2}$ & Open & Fokker 5-loop \\
$\zeta(3)$ in local $A$ at 7\PN+ & Open & Unknown \\
$\sigma_r=0$ dynamical consistency & Open & Check needed \\
$G_t$ from full action variation & Open & Target known \\
$c_{60}^{\rm rat}$ range & Open & $\approx14390\pm40$ \\
$\kappa$ from UV completion & Open & NJL $\sim10^{-4}$ \\
\bottomrule
\end{tabular}
\end{table}

%=====================================================================
\section{Conclusions}
%=====================================================================

Four new results extend the \QGD\ post-Newtonian programme:

\paragraph{1. Two-body CM field.}
$\sigma_{\rm total}^2 = 2u/\nu$ exactly at Newtonian order.
The \EOB\ canonical transformation corresponds to $\sigma_{\rm eff} = \sqrt{\nu}\,\sigma_{\rm total}$.

\paragraph{2. $B(u;\nu)=1$.}
In \QGD\ isotropic gauge, $g_{rr}=1$ for the two-body problem.
Identifying isotropic $r$ as the \EOB\ coordinate gives $B=1$.
This is a gauge statement but internally consistent and simplifies
the \EOB\ Hamiltonian to $H_{\rm eff} = \mu c^2\sqrt{A(1+p^2/\mu^2c^2)}$.

\paragraph{3. Physical interpretation of $\xi$.}
$\xi = 2275/656$ is the ratio of consecutive Hadamard momentum integrals;
the $u$-Pad\'e is the Borel resummation of the self-interaction series.
The analogy with the QCD Borel pole suggests the structure is robust.

\paragraph{4. Two of four missing BDG 2020 coefficients.}
\begin{equation*}
  a_{62}^{\pi^4} = \frac{184275}{131072}, \qquad
  a_{62}^{\pi^6} = -\frac{419225625}{85983232}.
\end{equation*}
These are the first explicit predictions of specific sub-coefficients
in the BDG 2020 unknown $\nu^3$ sector.
If confirmed (by higher-order \GSF\ numerics or completing the 5-loop
Fokker computation), they would represent the strongest test of the
\QGD\ $\beta$-tower to date.

\begin{thebibliography}{9}
\bibitem{matshaba2026}
R.~Matshaba,
\textit{Quantum Gravitational Dynamics}, UNISA (2026).
\bibitem{BD2025}
D.~Bini, T.~Damour, A.~Geralico,
arXiv:2507.08708 (2025).
\bibitem{BDG2020local}
D.~Bini, T.~Damour, A.~Geralico,
Phys.\ Rev.\ D \textbf{102}, 024061 (2020) [arXiv:2003.11891].
\bibitem{BDG2020nonlocal}
D.~Bini, T.~Damour, A.~Geralico,
Phys.\ Rev.\ D \textbf{102}, 084047 (2020) [arXiv:2007.11239].
\bibitem{BD2013}
D.~Bini, T.~Damour, Phys.\ Rev.\ D \textbf{87}, 121501(R) (2013).
\bibitem{DJS2001}
T.~Damour, P.~Jaranowski, G.~Sch\"afer, Phys.\ Lett.\ B \textbf{513}, 147 (2001).
\bibitem{DJS2015}
T.~Damour, P.~Jaranowski, G.~Sch\"afer, Phys.\ Rev.\ D \textbf{89}, 064058 (2014).
\bibitem{Buonanno1999}
A.~Buonanno, T.~Damour, Phys.\ Rev.\ D \textbf{59}, 084006 (1999).
\bibitem{Python2026}
R.~Matshaba, Claude,
\texttt{qgd\_master\_v4.py}, March 2026.
\end{thebibliography}

\appendix
\section{$\nu$-tower through $n=9$}
\begin{center}
\renewcommand{\arraystretch}{1.2}
\begin{tabular}{ccccc}
\toprule
$n$ & PN & $\nu^1$ ($k=0$) & $\nu^2$ ($k=1$) & $\nu^3$ ($k=2$) \\
\midrule
4 & 3PN & $-41/32$ & --- & --- \\
5 & 4PN & $+2275/512$ & $+20475/8192$ & --- \\
6 & 5PN & $-5175625/335872$ & $-46580625/5373952$ & $-419225625/85983232$ \\
7 & 6PN & $+11774546875/220332032$ & $\cdots$ & $\cdots$ \\
\bottomrule
\end{tabular}
\end{center}
All entries via $c_{n,k}=\beta^{2k}c_{n,0}$, $\beta^2=9/16$.
\end{document}
